\section{Frequency chain and ratio inversion}
\label{app:chain}

The observable is the out-of-loop beat between the transferred
\qty{1550}{nm} carrier and the local Yb-referenced comb. All
frequencies in the chain derive from the comb repetition rate
$f_{\mathrm{rep}} = \fRepMHz$\,MHz. The Sr clock's absolute frequency
$f_{\mathrm{Sr}}$ (\qty{698}{nm}) is related to the comb frequency fed
to the \qty{1397}{nm} branch by
\begin{equation}
  f_{\mathrm{comb}}^{1397} = \tfrac{1}{2}\,(1 + a)\,f_{\mathrm{Sr}},
\end{equation}
where $a$ collects the Sr systematic corrections (density, lattice,
second-order Zeeman and BBR shifts). The Yb clock's absolute frequency
$f_{\mathrm{Yb}}$ (\qty{578}{nm}) relates to the \qty{1156}{nm} comb
branch by
\begin{equation}
  f_{\mathrm{comb}}^{1156} = \tfrac{1}{2}\,(1 + b)\,f_{\mathrm{Yb}}.
\end{equation}
The beat is converted to fractional units with the transfer frequency
$F_{1550} = \fFifteenTHz$\,THz. The static gravitational-redshift
correction enters the ratio inversion as a multiplicative factor,
\begin{equation}
  \left(\mathrm{Yb/Sr}\right)_{\rm corr}
  = \left(\mathrm{Yb/Sr}\right)_{\rm raw} \times (1 + \delta_g),
  \qquad \delta_g = \deltaGV,
\end{equation}
derived from the levelled geopotential difference
$\dW = \dWVal \pm \dWUnc~\mathrm{m^2\,s^{-2}}$
(Section~\ref{sec:methods-levelling}).

\section{Data processing and segment definition}
\label{app:processing}

The counters sample the beat every second in phase-averaging mode.
Samples deviating by more than $10$\,Hz from the running median are
treated as jump points and removed; the retained data are further
trimmed at both ends by an endpoint screen that drops points beyond
$1\%$ of the peak-to-peak range of the record. The primary data set
consists of the longest jump-free record of each accepted segment
after this screening. Table~\ref{tab:segments} lists the per-segment
retained counts, ratio deviations and statistical uncertainties; the
category column marks the five segments whose stability fits are
model-limited. The excluded ninth recorded segment, its parameter
provenance and its effect on the combination are described in
Sections~\ref{sec:methods-campaign} and~\ref{sec:results-sensitivity}.

\begin{table}[t]
  \centering
  \caption{Per-segment data of the primary 16-segment set. Ratio
  deviations $y_i = R_i/R_1 - 1$ are relative to the first segment;
  $u(y_i)$ is the statistical uncertainty of the segment ratio;
  ``Fit'' is \emph{est.} for segments whose stability fit lies in the
  white-frequency-noise band and \emph{lim.} for model-limited fits.}
  \label{tab:segments}
  % Generated by tools/build_numbers.py -- do not edit by hand.
\begin{tabular}{rrrrl}
\toprule
Group & $n_{\rm valid}$ & $y_i$ ($10^{-18}$) & $u_i$ ($10^{-18}$) & Fit \\
\midrule
1 & 68\,009 & +0.00 & 1.113 & est. \\
2 & 29\,481 & -3.36 & 1.178 & lim. \\
3 & 42\,891 & +2.00 & 1.984 & est. \\
4 & 64\,280 & -4.61 & 0.951 & est. \\
5 & 19\,532 & -3.38 & 6.040 & est. \\
6 & 52\,140 & -1.15 & 0.513 & lim. \\
7 & 39\,526 & +5.31 & 1.173 & lim. \\
8 & 43\,565 & +0.97 & 4.509 & est. \\
10 & 57\,938 & -0.76 & 1.230 & est. \\
11 & 35\,824 & +2.61 & 2.301 & est. \\
12 & 40\,346 & +0.73 & 1.715 & lim. \\
13 & 152\,994 & -3.09 & 1.694 & est. \\
14 & 68\,618 & +1.97 & 1.377 & est. \\
15 & 57\,649 & -0.25 & 1.293 & est. \\
16 & 147\,640 & -0.77 & 0.607 & lim. \\
17 & 65\,983 & -0.89 & 2.142 & est. \\
\bottomrule
\end{tabular}

\end{table}

\section{Statistical combination models}
\label{app:stats}

\paragraph{Fixed-effect WLS.} Under the assumption that all segments
estimate one common ratio with uncertainties $u_i$,
\[
  \mu_{\rm WLS} = \frac{\sum_i w_i R_i}{\sum_i w_i}, \quad
  w_i = 1/u_i^2, \quad
  u_{\rm WLS} = \left(\sum_i w_i\right)^{-1/2},
\]
with $\chi^2 = \sum_i (R_i - \mu_{\rm WLS})^2 / u_i^2$ and
$\chi^2_{\rm red} = \chi^2/(N-1)$. For the primary data set,
$\chi^2_{\rm red} = \chiRed$ ($N = \nSeg$,
$\dofV$ degrees of freedom, $p = \pChi$), so the fixed-effect
uncertainty under-describes the observed scatter.

\paragraph{Birge scaling.} The Birge ratio
$B = \sqrt{\chi^2_{\rm red}} = \birgeRatio$ scales all segment
uncertainties by a common factor, $u_i \to B u_i$; it assumes the
relative sizes of the $u_i$ are correct but their overall scale is
not.

\paragraph{Mandel--Paule.} An excess scatter $\xi$ common to all
segments is added in quadrature, $u_{i,\rm eff}^2 = u_i^2 + \xi^2$,
and $\xi$ is solved from $\chi^2_{\rm red}(\xi) = 1$. For the primary
data set, $\xi = \xiMpE \times 10^{-18}$.

\paragraph{Bayesian random effect.} The same observation model,
$R_i \sim \mathcal{N}(\mu, u_i^2 + \xi^2)$, is treated
hierarchically: $\mu$ and $\xi$ are both inferred, with flat priors
on $\mu$ and on $\log\xi$, and the joint posterior
$p(\mu, \xi \mid \{R_i, u_i\})$ is marginalised numerically over a
40-decade log-$\xi$ grid. The reported central value is the posterior
mean of $\mu$ and the statistical uncertainty is its posterior
standard deviation; for the primary data set,
$\xi = \xiBayesE \times 10^{-18}$. The excess scatter quantified here
also absorbs the correlated component of the time-varying
gravitational-redshift term (Section~\ref{sec:methods-budget}), which
is carried separately as an uncertainty item rather than corrected.

\section{Levelling reduction}
\label{app:levelling}

The orthometric height difference between benchmarks follows from the
levelled normal heights after the normal-to-orthometric conversion,
with gravity observed at every levelling point. The geopotential
difference is integrated segment by segment,
$\dW_{AB} = -\sum_i g_i \Delta H_i$, and combined with the two local
offsets between benchmark and clock platform (obtained by
trigonometric levelling) and the platform-to-transition offsets from
the clock design parameters. The resulting atomic-transition
geopotential difference is $\dW = \dWVal \pm \dWUnc~\mathrm{m^2\,s^{-2}}$,
equivalent to an optical fractional shift of $\deltaGV$.

\section{Uncertainty budget details}
\label{app:budget}

The components of Table~\ref{tab:budget} are obtained as follows.
The statistical term is the posterior standard deviation of the
Bayesian random-effect model (Appendix~\ref{app:stats}). The Sr and
Yb terms are the published systematic uncertainties of the two
clocks~\citep{jia2026sr, zhu2026yb}; for the Yb clock the published
evaluation is adopted. The static-potential term is
$u(\dW)/c^2$ from the levelling uncertainty
$\dWUnc~\mathrm{m^2\,s^{-2}}$, giving
$\uStaticE \times 10^{-18}$. The fibre-link and comb terms are
experiment-side upper bounds of $\uSubsystemBound$ each, carried at
\uLinkE $\times 10^{-18}$ and $\uCombE \times 10^{-18}$ to preserve the
experiment-side margin. The time-varying-redshift term is scaled from
the spread of the segment-averaged gravitational shift across the
campaign and is carried at $\uTidalE \times 10^{-18}$; it is not
applied as a correction. The total standard uncertainty,
$\uTotalE \times 10^{-18}$, is the quadrature sum of the seven
components.
